\documentclass{article}
\usepackage[utf8]{inputenc}
% %\input{GrandMacros.tex}
% %\input{Macro.tex}
\usepackage{amsthm}
\usepackage{xcolor}
\newtheorem{remark}{Remark}
\newtheorem{lemma}{Lemma}
\newtheorem{theorem}{Theorem}
\newtheorem{proposition}{Proposition}
\newtheorem{corollary}{Corollary}
\usepackage{graphicx,verbatim,array,multicol, subfigure, color, lscape, mathrsfs, dsfont,url}
%\usepackage{psfrag, amsmath, amsfonts, epsfig, fancybox, setspace,soul, cancel,biometri}
\usepackage{psfrag, amsmath, amsfonts, epsfig, fancybox, setspace, soul, amsthm, longtable, threeparttable, threeparttablex, makecell, xcolor, pdflscape, tabu, booktabs, natbib, amssymb}
\usepackage[normalem]{ulem}
\usepackage{cancel}
\usepackage[utf8]{inputenc}
\parindent 16pt
\include{common-defs}
\newcommand{\PP}{\mathbb{P}}

\input{typography-preamble}

% \def\cor{\text{cor}}
% \def\cF{\cF}

\title{Surrogate inference via random probability distributions: evaluation of surrogate quality in future studies}


\begin{document}
\maketitle

\section{Introduction}

A central question in the evaluation of surrogate markers is whether the surrogate can replace the gold-standard outcome in future trials or in new populations. To inform this decision, a range of competing approaches have been proposed \cite{joffe2009related}. Each approach comes with assumptions that are difficult to assess in practice (e.g., no unmeasured confounding, cross-world independence) or with data requirements that limit implementation. For example, meta-analytic validation approaches \cite{buyse2000validation} rely on the existence of multiple studies that estimate both the treatment effect on the outcome and the treatment effect on the surrogate. In many instances, such studies are not available. Even when multiple studies are available, it may not be known a priori when enough studies---capturing relevant treatment effect heterogeneity---have been accumulated. 

Other strategies rely on untestable assumptions on the distribution of counterfactual quantities without which point identification is not possible. 
\begin{itemize}
    \item Principal stratification can require strong assumptions on the joint distribution of $\{Y(1), Y(0)\}$ and $\{S(1), S(0)\}$, e.g., conditional independence of principal strata given covariates, which are often regarded as strong \cite{frangakis2002principal, kim2024bayesian}.
    \item Other approaches to principal stratification make restrictive assumptions on the nature of the counterfactual $S(0)$ (e.g., that it is constant) as well as requiring near-perfect prediction performance of $S(1)$ \cite{elliott2023surrogate}. 
    \item Causal mediation methods require no unobserved confounding between the surrogate and outcome \cite{vanderweele2015mediation}, which can be difficult to ensure in practice. 
    \item Related methods based on the proportion of treatment effect (PTE) do not require such unobserved confounding assumptions, but they do require restrictive monotonicity assumptions to avoid the surrogate paradox \cite{chen2003pte}.
\end{itemize}

A central motivation for surrogate evaluation is use in future trials or in new populations. A surrogate is not needed in the current study because the gold-standard outcome is observed there; interest often lies in a future study on a new population (with unknown characteristics), where the gold standard may be unobserved. The target of inference is therefore not the current study alone, but the behavior of the surrogate--outcome relationship across such future populations.

In this paper we formalize inference about surrogate quality in future studies. We treat the current study's distribution as one draw from a super-population of study-specific distributions, and we define a random probability measure $\cF$ over such distributions. We specify the law of $\cF$ via a perturbation model (future studies as mixtures of the current distribution and an innovation), establish representation and continuity results, and provide a plug-in estimator with asymptotic normality. We also give a procedure to report a distance $\varepsilon$ within which the surrogate remains ``good'' (e.g., correlation of treatment effects above a threshold). The development is primarily frequentist; a Bayesian implementation (e.g., posterior over $\phi(\cF)$) can be obtained by placing a prior on the innovation and is left for discussion.

Table~\ref{tab:comparison} summarizes how this approach compares to meta-analytic validation, principal stratification, PTE-based methods, and causal mediation in terms of inferential target, main assumptions, and data requirements. Our approach targets functionals of the \emph{distribution} of studies (e.g., correlation of treatment effects across studies), requires only a single current study plus a specification of how future studies may differ (the perturbation model), and avoids cross-world and unmeasured-confounding assumptions needed for principal stratification and mediation.

\begin{table}[htbp]
\centering
\caption{Comparison of surrogate evaluation approaches.}
\label{tab:comparison}
\begin{tabular}{llll}
\toprule
Approach & Inferential target & Main assumptions & Data \\
\midrule
Meta-analytic \cite{buyse2000validation} & Trial-level association & Multiple trials; homogeneity & Multiple studies \\
Principal stratification \cite{frangakis2002principal} & Principal causal effects & Principal strata structure; often strong independence & Single/multiple \\
PTE \cite{chen2003pte} & Proportion of effect explained & Monotonicity (to avoid surrogate paradox) & Single study \\
Mediation \cite{vanderweele2015mediation} & Natural direct/indirect effects & No unmeasured confounding (M--Y; A--M--Y) & Single study \\
This approach & Functionals of distribution of studies & Perturbation model; finite support (estimation) & Single study \\
\bottomrule
\end{tabular}
\end{table}

\section{Setting}\label{sec:setting}
Suppose we observe $n$ iid realizations of $\bO = (\bX, A, S, Y) \sim \PP_0$ in the current study, where $A \in \{0, 1\}$ is a binary treatment, $Y \in \cY \subset \R$ is the outcome of interest, $S \in \cS \subset \R$ is a surrogate marker that may replace $Y$ in a future study, and $\bX \in \cX \subset \R^p$ are pre-treatment covariates ($\bX$ may be omitted in a randomized trial). The distribution $\PP_0$ is the law of the observables. We measure treatment effects on the surrogate and the outcome via
\begin{align*}
    &\Delta_Y = \E_{\PP_0}\{Y(1) - Y(0)\}\\
    &\Delta_S = \E_{\PP_0}\{S(1) - S(0)\}
\end{align*}
This all describes (the population of) the current study. Under suitable identification assumptions, $\Delta_S$ and $\Delta_Y$ are identified and can be estimated by standard methods (e.g., difference in means or influence-function-based estimators). However, for the purposes of determining surrogate quality, we wish to know how the surrogate would perform in \emph{future} studies. Consider a future study with distribution $\PP^*$; we would observe $\Delta^*_S = \E_{\PP^*}\{S(1) - S(0)\}$ but not $Y$, and we wish to know what $\Delta^*_S$ implies about $\Delta^*_Y = \E_{\PP^*}\{Y(1) - Y(0)\}$. 

To formalize this, we must begin to characterize a class of future studies to which we might wish to generalize. We treat the study-specific distribution as a draw from a random probability measure $\cF$, so that both the current distribution $\PP_0$ and any future study distribution $\Q$ are (conceptually) draws from $\cF$.

Suppose that we have a random probability measure $\cF$ such that $\PP_0 \sim \cF$ and any future study $\Q \sim \cF$. Our goal is to make inference about functionals $\phi(\cF)$ of the distribution of studies using the observed data from the current study. Examples of functionals that may be of interest include the correlation between the two treatment effects $\phi(\cF) = \frac{\Cov_{\cF}\{\Delta_S(\Q), \Delta_Y(\Q)\}}{\sqrt{\Var_\cF\{\Delta_S(\Q)\}\Var_\cF\{\Delta_Y(\Q)\}}}$; the probability that a non-small treatment effect on the surrogate implies a non-small treatment effect on the outcome $\phi(\cF; \epsilon_S, \epsilon_Y) = \E_\cF\left[ \cI\{ \Delta_S(\Q) > \epsilon_S, \Delta_Y(\Q) > \epsilon_Y\} \right] / \E_\cF\left[ \cI\{ \Delta_S(\Q) > \epsilon_S\} \right]$; or the expected outcome treatment effect for a given surrogate treatment effect $\phi(\cF; \delta) = \E_\cF\left\{\Delta_Y(\Q) | \Delta_S(\Q) = \delta\right\}$. Our goal is to estimate these functionals, thereby characterizing how the surrogate--outcome relationship may behave in future studies.

\subsection{Defining the random probability measure $\cF$}

Let $\Omega = \cX \times \{0,1\} \times \cS \times \cY$ denote the observed data space, and let $\cM_1(\Omega)$ denote the space of all probability measures on $\Omega$. The current study distribution $\PP_0$ is one element of $\cM_1(\Omega)$; future study distributions $\Q$ are also elements of this space.

We operate under the assumption that future studies are likely to bear \emph{some} resemblance to the current study -- if the current study is to be informative of the future study, they should have some similarity. We encode this future similarity by modelling future studies via a perturbation of $\PP_0$. For a fixed perturbation level $\lambda \in [0,1]$, define the $\lambda$-perturbation family
\begin{align*}
    \cF_\lambda(\PP_0) = \{\Q \in \cM_1(\Omega) : \Q = (1-\lambda)\PP_0 + \lambda\tilde{\PP} \text{ for some } \tilde{\PP} \in \cM_1(\Omega)\}.
\end{align*}
This is the set of all convex combinations of $\PP_0$ with weight $(1-\lambda)$, where the innovation $\tilde{\PP}$ can be any probability measure. The parameter $\lambda$ controls how far future studies can deviate from $\PP_0$: when $\lambda = 0$, only $\Q = \PP_0$ is possible; when $\lambda = 1$, any distribution is reachable.

To make $\cF_\lambda$ a random probability measure, we endow it with a probability distribution by specifying a law $\mu$ on the innovation $\tilde{\PP}$. Given such a distribution $\mu$ on $\cM_1(\Omega)$, we write
\begin{align*}
    \Q \sim \cF_\lambda \quad\Leftrightarrow\quad \Q = (1-\lambda)\PP_0 + \lambda\tilde{\PP}, \quad \tilde{\PP} \sim \mu.
\end{align*}

The perturbation form is general in the sense that every future-study distribution admits the representation above. For any probability measure $\Q \in \cM_1(\Omega)$, there exists $\lambda \in [0,1]$ such that $\Q = (1-\lambda)\PP_0 + \lambda\tilde{\PP}$ for some $\tilde{\PP} \in \cM_1(\Omega)$. However, once we endow $\cF_\lambda$ with a probability distribution via $\mu$, the \emph{support} of the random measure $\cF_\lambda$ is restricted to those $\Q$ where $\tilde{\PP}$ lies in the support of $\mu$. Thus every $\Q$ is representable, but the \emph{probabilistic support} under a given $\mu$ depends on the support of $\mu$: if $\mu$ has full support on $\cM_1(\Omega)$, the random measure can reach any distribution; if $\mu$ is concentrated (e.g., near $\PP_0$), then only a subset of $\cM_1(\Omega)$ receives positive probability. The choice of $\mu$ therefore determines which directions of deviation from $\PP_0$ are probabilistically plausible, not merely representable.

The parameter $\lambda$ controls the distance between future studies and the current study. For any $\Q = (1-\lambda)\PP_0 + \lambda\tilde{\PP}$, we have
\begin{align}\label{eq:distance-bound}
    d_{\mathrm{TV}}(\Q, \PP_0) \leq \lambda \, d_{\mathrm{TV}}(\tilde{\PP}, \PP_0) \leq \lambda
\end{align}
where $d_{\mathrm{TV}}(\PP, \Q) = \sup_{A \in \Omega} \left|\PP(A) - \Q(A)\right|$ is the total variation distance and $d_{\mathrm{TV}}(\PP, \Q) \leq 1$ for any two probability measures $\PP, \Q \in \cM_1(\Omega)$. Thus, fixing $\lambda$ restricts attention to distributions within total variation distance $\lambda$ of $\PP_0$: $\cF_\lambda(\PP_0) = \{\Q \in \cM_1(\Omega) : d_{\mathrm{TV}}(\Q, \PP_0) \leq \lambda\}$. When $\lambda$ is close to 0, the distribution on $\cF_\lambda$ induced by $\mu$ concentrates tightly around $\PP_0$; when $\lambda = 0$, $\Q = \PP_0$ with probability 1. As $\lambda$ approaches 1, the distribution of $\Q$ approaches that of $\tilde{\PP}$ and may differ substantially from $\PP_0$. For intermediate $\lambda$, $\Q$ interpolates between $\PP_0$ and $\tilde{\PP}$. Thus, the choice of $\lambda$ and the distribution $\mu$ on $\tilde{\PP}$ jointly encode which future studies are plausible and how far they may deviate from the current study.

% \subsubsection{Choosing the innovation distribution $\mu$}
%
% The interpretation of surrogate quality depends critically on the choice of $\mu$, not just $\lambda$. If $\mu$ concentrates mass near $\PP_0$ (e.g., $\mu = \text{Dirichlet}(\alpha \bp_0)$ with large $\alpha$), then finding $\phi(\cF_\lambda) \geq c$ for large $\lambda$ provides only weak evidence---it merely says the surrogate performs well for future studies that closely resemble $\PP_0$. A more stringent test uses a non-informative $\mu$ that treats all innovations $\tilde{\PP} \in \cM_1(\Omega)$ equally, providing no favored direction of deviation from $\PP_0$.
%
% For computational tractability and to define a natural non-informative prior, we assume finite support:

\noindent\textbf{Assumption (Finite support).}\label{ass:finite}
The outcome space $\Omega$ has finite support. Equivalently, the observed data naturally partitions into a finite number of cells (e.g., based on observed covariate values, treatment, surrogate, and outcome levels).

This assumption is reasonable in practice: with $n$ observations, empirical distributions are inherently discrete; many applications involve categorical or discretized outcomes; and finite support makes inference tractable. Under finite support, $\cM_1(\Omega)$ is a $(k-1)$-simplex (where $k = |\Omega|$).

\subsubsection{Treatment effect heterogeneity as the fundamental object}

While the finite support assumption allows for tractable computation, it is important to recognize that the \emph{fundamental source of variation} across studies is heterogeneity in treatment effects, not arbitrary partitions of the data space. For continuous covariates $\bX$, treatment effects $\tau_S(\bX) = S(1) - S(0)$ and $\tau_Y(\bX) = Y(1) - Y(0)$ are continuous functions. Future studies differ from the current study primarily through different distributions of $\bX$, which induce different distributions of treatment effects.

The finite support assumption operationalizes this by discretizing $\bX$ into $J$ bins and treating each bin as a cell. The choice of $J$ represents an approximation: as $J \to \infty$ (with $J = o(n)$), the discretized model approaches the continuous treatment effect function. In practice, $J$ should be chosen large enough to capture meaningful heterogeneity but small enough to avoid over-discretization (which increases estimation noise).

This perspective clarifies that there are no ``true'' types---only varying treatment effects. Discretization schemes that align with treatment effect heterogeneity (e.g., using random forests to partition based on estimated $\tau(\bX)$) will provide better approximations than arbitrary covariate-based discretizations. Different discretization schemes explore different aspects of the treatment effect distribution; using multiple schemes and taking their minimum (ensemble approach) provides a more comprehensive assessment of worst-case surrogate performance.

% and we can define a canonical non-informative prior.
%
% Under finite support, the natural choice for a non-informative $\mu$ is the uniform Dirichlet:
% \begin{align}
%     \mu = \text{Dirichlet}(1, \ldots, 1).
% \end{align}
% This distribution is uniform on the simplex $\cM_1(\Omega)$, treating all probability vectors equally. It is the maximum entropy prior given no constraints, providing the most conservative (stringent) assessment of surrogate quality. When $\phi(\cF_\lambda) \geq c$ under this choice of $\mu$, the surrogate performs well across a broad, evenly-distributed set of future studies within distance $\lambda$ of $\PP_0$.
%
% \subsubsection{How similar must a future study be for the surrogate to remain good?}
%
% A central question is: \emph{how different can a future study be from the current study while the surrogate remains good?} We formalize this by asking for what values of $\lambda$ the functional $\phi(\cF_\lambda)$ stays above a desired threshold. For example, if we define ``good surrogate'' as correlation of treatment effects above 0.8, we want to find the largest $\lambda$ for which $\phi(\cF_\lambda) \geq 0.8$. This $\lambda$ tells us how far a future study can deviate from the current study (in total variation distance) while maintaining adequate surrogate quality.
%
% To make this operational, we need to understand how functionals $\phi(\cF_\lambda)$ behave as $\lambda$ varies. The key insight is that the treatment effects $\Delta_S(\Q)$ and $\Delta_Y(\Q)$ depend continuously on $\Q$ (in the sense that small changes in $\Q$ produce small changes in the treatment effects), and thus functionals of the treatment effects also vary smoothly with $\lambda$. The following assumptions and result formalize this.
%
% \noindent\textbf{Assumption 3 (Bounded outcomes).}
% $S \in [s_{\min}, s_{\max}]$ and $Y \in [y_{\min}, y_{\max}]$ are bounded.
%
% Under this assumption, treatment effects are Lipschitz continuous in $\Q$. For any $\Q, \Q'$,
% \begin{align}\label{eq:lipschitz}
%     |\Delta_S(\Q) - \Delta_S(\Q')| &\leq 2(s_{\max} - s_{\min}) \, d_{\mathrm{TV}}(\Q, \Q')\\
%     |\Delta_Y(\Q) - \Delta_Y(\Q')| &\leq 2(y_{\max} - y_{\min}) \, d_{\mathrm{TV}}(\Q, \Q'). \nonumber
% \end{align}
% This follows because expectations of bounded functions are Lipschitz in total variation: $|\E_{\Q} f - \E_{\Q'} f| \leq (\sup f - \inf f) \, d_{\mathrm{TV}}(\Q, \Q')$. Under SUTVA and consistency, $S(1) - S(0)$ ranges over $[-(s_{\max} - s_{\min}), (s_{\max} - s_{\min})]$, yielding the constant $2(s_{\max} - s_{\min})$, and similarly for $Y$.
%
% \begin{theorem}[Surrogate quality under perturbations]\label{thm:epsilon}
% Suppose ``good surrogate'' is defined by $\phi(\cF_\lambda) \geq c$ for some threshold $c$ (e.g., $\phi$ is the correlation of treatment effects and $c = 0.8$). Under Assumption 3, the functional $\phi(\cF_\lambda)$ varies continuously with $\lambda$. If $\phi(\cF_\lambda) \geq c$ for some $\lambda$, then any future study $\Q$ with $d_{\mathrm{TV}}(\Q, \PP_0) \leq \lambda$ will have $|\phi(\Q) - \phi(\cF_\lambda)| \leq L \lambda$ for some constant $L$ depending on $\phi$ and regularity conditions (e.g., non-degenerate variances for correlation).
% \end{theorem}
% \begin{proof}
% By \eqref{eq:lipschitz}, the treatment effects $(\Delta_S(\Q), \Delta_Y(\Q))$ are Lipschitz continuous in $\Q$ with respect to $d_{\mathrm{TV}}$. The functionals $\phi$ of interest (correlation, probability of concordance, conditional expectation) are continuous---and under bounded $S$, $Y$ and non-degenerate variances, Lipschitz---as functions of the bivariate distribution of treatment effects. Since $\Q \in \cF_\lambda(\PP_0)$ implies $d_{\mathrm{TV}}(\Q, \PP_0) \leq \lambda$ by \eqref{eq:distance-bound}, the distribution of $(\Delta_S(\Q), \Delta_Y(\Q))$ varies continuously with $\lambda$, and composing the Lipschitz maps yields continuity of $\phi(\cF_\lambda)$ in $\lambda$.
% \end{proof}
%
% \begin{remark}[Operationalizing the threshold]\label{remark:operational}
% In practice, we estimate $\phi(\cF_\lambda)$ for a range of $\lambda$ values (e.g., $\lambda \in \{0.05, 0.1, 0.15, \ldots, 0.5\}$) by generating future studies $\Q = (1-\lambda)\hat{\PP}_n + \lambda\tilde{\PP}$ with $\tilde{\PP} \sim \mu$ (e.g., $\mu = \text{Dirichlet}(1, \ldots, 1)$), computing $(\Delta_S(\Q), \Delta_Y(\Q))$ for each, and calculating the empirical correlation (or other functional). The largest $\lambda$ for which $\hat{\phi}(\cF_\lambda) \geq c$ provides an assessment: \emph{when future studies are drawn from $\cF_\lambda$ (within distance $\lambda$ and distributed according to $\mu$), the surrogate performs well on average.} This operationalizes ``how different future studies can be while maintaining adequate surrogate quality.''
%
% Importantly, $\phi(\cF_\lambda) \geq c$ is a statement about the \emph{average} surrogate performance over the distribution $\mu$, not a guarantee for every individual study within distance $\lambda$. With uniform $\mu = \text{Dirichlet}(1, \ldots, 1)$, this average represents a stringent, conservative test.
% \end{remark}
%
% \begin{remark}[Sensitivity to $\mu$]\label{remark:mu-sensitivity}
% We recommend reporting $\hat{\phi}(\cF_\lambda)$ under multiple choices of $\mu$ to assess sensitivity. For example:
% \begin{itemize}
%     \item $\mu_{\text{uniform}} = \text{Dirichlet}(1, \ldots, 1)$ (non-informative, most stringent)
%     \item $\mu_{\text{moderate}} = \text{Dirichlet}(\alpha \bp_0)$ with moderate $\alpha$ (e.g., $\alpha = 5$)
%     \item $\mu_{\text{concentrated}} = \text{Dirichlet}(\alpha \bp_0)$ with large $\alpha$ (e.g., $\alpha = 50$)
% \end{itemize}
% If $\hat{\phi}(\cF_\lambda) \geq c$ for large $\lambda$ even under $\mu_{\text{uniform}}$, this provides strong evidence of robust surrogate quality. If this holds only under $\mu_{\text{concentrated}}$, the evidence is weaker (surrogate is good only for futures resembling $\PP_0$).
% \end{remark}

\subsection{Estimation}\label{sec:estimation}

\subsubsection{Setup for fixed $\lambda$}

For a fixed value $\lambda \in [0,1]$, we define the perturbation family
$\cF_\lambda(\PP_0)$ as in Section~\ref{sec:setting}. The estimand is the functional
$\phi(\cF_\lambda)$ evaluated at this fixed $\lambda$. We treat $\lambda$
as a design parameter—analogous to choosing a significance level—that
controls how far future studies may deviate from $\PP_0$.

In this section, we develop asymptotic theory for estimating $\phi(\cF_\lambda)$
for a single fixed $\lambda$. Section~\ref{sec:grid} addresses the practical
question of searching over multiple $\lambda$ values to find the largest
$\lambda^*$ for which $\phi(\cF_{\lambda^*}) \geq c$.

\subsubsection{Regularity conditions}
We assume the following.

\noindent\textbf{Assumption 1 (Bounded outcomes).}\label{ass:bounded}
$S \in [s_{\min}, s_{\max}]$ and $Y \in [y_{\min}, y_{\max}]$ almost surely
under $\PP_0$ and all $\Q \in \cF_\lambda(\PP_0)$, with
$0 < s_{\max} - s_{\min} < \infty$ and $0 < y_{\max} - y_{\min} < \infty$.

\noindent\textbf{Assumption 2 (Identification and estimability).}\label{ass:identification}
The treatment effects $\Delta_S(\Q)$ and $\Delta_Y(\Q)$ are identified from
observables under standard causal identification conditions (e.g., randomization,
unconfoundedness, or instrumental variables), and admit $\sqrt{n}$-consistent
estimators with influence functions satisfying $\E[\psi^2] < \infty$.

\noindent\textbf{Assumption 3 (Functional smoothness).}\label{ass:smoothness}
For the correlation functional, the function
$H(\delta_S, \delta_Y) = \E_\mu[\rho(\Delta_S(\Q), \Delta_Y(\Q))]$
(where the expectation is over $\tilde{\PP} \sim \mu$ in
$\Q = (1-\lambda)\PP_0 + \lambda\tilde{\PP}$, and $\rho$ denotes correlation)
is continuously differentiable in $(\delta_S, \delta_Y)$ with bounded gradient
$\|\nabla H\| < \infty$ in a neighborhood of $(\Delta_S(\PP_0), \Delta_Y(\PP_0))$.

For the probability and conditional mean functionals, $H$ is Lipschitz continuous
with constant $L_H < \infty$.

\noindent\textbf{Assumption 4 (Non-degeneracy).}\label{ass:nondegeneracy}
For the correlation functional, there exists $\epsilon > 0$ such that
\begin{align*}
    \inf_{\Q \in \cF_\lambda(\PP_0)} \Var_\mu[\Delta_S(\Q)] &\geq \epsilon\\
    \inf_{\Q \in \cF_\lambda(\PP_0)} \Var_\mu[\Delta_Y(\Q)] &\geq \epsilon
\end{align*}
where the variance is taken over the innovation distribution $\mu$
(not within a single study).

\begin{remark}[On finite support]
Assumption~\ref{ass:finite} (finite support, stated in Section~\ref{sec:setting}) is restrictive
and primarily computational. It enables tractable simulation and a natural uniform
prior $\mu = \text{Dirichlet}(1,\ldots,1)$. In practice, continuous outcomes can be
discretized into a fine grid (trading bias for computational feasibility), or
replaced by kernel-based methods (left for future work). The asymptotic theory
extends to growing support $k=k_n$ under additional rate conditions, but we do not
pursue that here.
\end{remark}

Under Assumption~\ref{ass:finite}, the empirical distribution $\hat{\PP}_n$ is the vector of empirical proportions $\hat{\bp}_n$; $\sqrt{n}(\hat{\bp}_n - \bp_0)$ is asymptotically normal (multinomial).

\subsubsection{Asymptotic theory}

\paragraph{The plug-in estimator.}
For fixed $\lambda$, define the plug-in estimator
\[
    \hat{\phi}_n(\lambda) = \phi(\cF_\lambda(\hat{\PP}_n))
\]
where $\cF_\lambda(\hat{\PP}_n)$ is the perturbation family centered at the
empirical distribution $\hat{\PP}_n$ instead of $\PP_0$. Operationally, we:
\begin{enumerate}
    \item Generate $M$ innovations $\tilde{\PP}_m \sim \mu$ (e.g., Dirichlet)
    \item Form $\Q_m = (1-\lambda)\hat{\PP}_n + \lambda\tilde{\PP}_m$
    \item Compute $(\hat{\Delta}_S(\Q_m), \hat{\Delta}_Y(\Q_m))$ for each $\Q_m$
    \item Compute $\hat{\phi}_n(\lambda)$ as the sample correlation (or other functional)
          of the $M$ pairs
\end{enumerate}

\paragraph{Variance decomposition.}
The variance of $\hat{\phi}_n(\lambda)$ arises from three sources:
\begin{enumerate}
    \item \textbf{Estimation error in $\hat{\PP}_n$:} The empirical distribution
          differs from $\PP_0$ by $O_p(n^{-1/2})$.
    \item \textbf{Monte Carlo error in simulating $\cF_\lambda$:} We use $M$ draws
          from $\mu$ to approximate $\E_\mu[\cdot]$, contributing error $O_p(M^{-1/2})$.
    \item \textbf{Treatment effect estimation error:} Within each $\Q_m$, we estimate
          $\Delta_S(\Q_m)$ and $\Delta_Y(\Q_m)$ from finite samples (if sampling
          from $\Q_m$ rather than using mixture weights directly).
\end{enumerate}

Under Assumptions~\ref{ass:bounded}--\ref{ass:smoothness}, the dominant term is
(1) when $M \to \infty$.

\paragraph{Asymptotic normality.}
Because $\Delta_S(\Q) = (1-\lambda)\Delta_S(\PP_0) + \lambda\Delta_S(\tilde{\PP})$ and similarly for $\Delta_Y$, the distribution of $(\Delta_S(\Q), \Delta_Y(\Q))$ under $\cF_\lambda$ depends on $\PP_0$ only through $(\Delta_S(\PP_0), \Delta_Y(\PP_0))$. Thus $\phi(\cF_\lambda)$ can be written as $\phi(\cF_\lambda) = H(\Delta_S(\PP_0), \Delta_Y(\PP_0))$ for some function $H$ determined by the fixed $\lambda$ and the distribution $\mu$.

\begin{proposition}\label{prop:asymptotic}
Under Assumptions~\ref{ass:bounded}--\ref{ass:nondegeneracy}, for fixed $\lambda$
and $M \to \infty$,
\[
    \sqrt{n}(\hat{\phi}_n(\lambda) - \phi(\cF_\lambda))
    \xrightarrow{d} N(0, \sigma^2(\lambda))
\]
where
\[
    \sigma^2(\lambda) = (\nabla H)^\top V(\lambda) (\nabla H)
\]
and:
\begin{itemize}
    \item $H$ is the function such that $\phi(\cF_\lambda) = H(\Delta_S(\PP_0), \Delta_Y(\PP_0))$
    \item $\nabla H$ is the gradient of $H$ evaluated at $(\Delta_S(\PP_0), \Delta_Y(\PP_0))$
    \item $V(\lambda)$ is the $2 \times 2$ asymptotic covariance matrix of
          $\sqrt{n}(\hat{\Delta}_S(\hat{\PP}_n) - \Delta_S(\PP_0),
                    \hat{\Delta}_Y(\hat{\PP}_n) - \Delta_Y(\PP_0))^\top$
\end{itemize}
\end{proposition}

\begin{proof}
By Assumption~\ref{ass:identification}, the treatment effect estimators satisfy
\[
    \sqrt{n}\begin{pmatrix} \hat{\Delta}_S(\hat{\PP}_n) - \Delta_S(\PP_0) \\
                            \hat{\Delta}_Y(\hat{\PP}_n) - \Delta_Y(\PP_0)
    \end{pmatrix}
    \xrightarrow{d} N(0, V(\lambda))
\]
where $V(\lambda)$ is the influence-function-based variance (standard in
semiparametric theory).

By Assumption~\ref{ass:smoothness}, the delta method applies:
\[
    \sqrt{n}(\hat{\phi}_n(\lambda) - \phi(\cF_\lambda))
    = (\nabla H)^\top \sqrt{n}\begin{pmatrix} \hat{\Delta}_S - \Delta_S \\
                                              \hat{\Delta}_Y - \Delta_Y
                      \end{pmatrix} + o_p(1)
    \xrightarrow{d} N(0, (\nabla H)^\top V(\lambda) (\nabla H))
\]
\end{proof}

\paragraph{Bootstrap inference.}
Direct computation of $V(\lambda)$ requires influence functions and may be
unstable in small samples. We recommend the \textit{nested bootstrap}:
\begin{enumerate}
    \item For $b = 1, \ldots, B$ bootstrap replicates:
          \begin{enumerate}
              \item Resample the data with replacement to get $\hat{\PP}_n^{(b)}$
              \item Generate $M$ future studies $\Q_m^{(b)} = (1-\lambda)\hat{\PP}_n^{(b)} + \lambda\tilde{\PP}_m$
              \item Compute $\hat{\phi}_n^{(b)}(\lambda)$
          \end{enumerate}
    \item Use the bootstrap distribution $\{\hat{\phi}_n^{(b)}(\lambda)\}_{b=1}^B$
          for inference (percentile intervals or standard errors)
\end{enumerate}

\begin{remark}[Bootstrap validity]
The nested bootstrap is valid under the conditions of Proposition~\ref{prop:asymptotic}.
The two-layer randomness (resampling data + drawing innovations) correctly
mimics the two-layer structure of the plug-in estimator. Standard bootstrap
theory applies.
\end{remark}

\section{Minimax bounds: eliminating dependence on $\mu$}\label{sec:minimax}

The inference framework in Section~\ref{sec:estimation} depends on a choice of innovation distribution $\mu$. The functional $\phi(\cF_\lambda)$ represents the \emph{expected} surrogate quality over $\mu$, which depends on our belief about which future studies are plausible. If $\mu$ concentrates near $\PP_0$, then $\phi(\cF_\lambda) \geq c$ provides only weak evidence (the surrogate is good for futures resembling the current study). A more stringent test uses a broad $\mu$ (e.g., uniform Dirichlet), but this still requires specifying $\mu$ correctly.

We now eliminate this $\mu$-dependence by providing \emph{worst-case bounds} that hold for \emph{all} innovation distributions in a plausible class $\cM$. This minimax approach guarantees coverage without requiring correct specification of $\mu$.

\subsection{Minimax inference framework}

Fix $\lambda \in [0,1]$ and define a class of innovation distributions
\[
    \cM = \{\mu : \text{$\mu$ is a probability measure on $\cM_1(\Omega)$ satisfying certain constraints}\}.
\]
For example, $\cM$ might be all Dirichlet distributions with concentration parameter $\alpha \in [\alpha_{\min}, \alpha_{\max}]$, or all distributions concentrated on individual observations (point masses), or some combination.

For each $\mu \in \cM$, we can compute $\phi(\cF_\lambda^\mu)$ where $\cF_\lambda^\mu$ is the random probability measure with innovation distribution $\mu$. The minimax bounds are:
\begin{align}
    \phi_*(\lambda) &= \inf_{\mu \in \cM} \phi(\cF_\lambda^\mu) \label{eq:phi-star-lower}\\
    \phi^*(\lambda) &= \sup_{\mu \in \cM} \phi(\cF_\lambda^\mu) \label{eq:phi-star-upper}
\end{align}

These bounds satisfy the following guarantee:

\begin{theorem}[Minimax coverage guarantee]\label{thm:minimax}
For any fixed $\lambda$ and class $\cM$, the interval $[\phi_*(\lambda), \phi^*(\lambda)]$ contains $\phi(\cF_\lambda^\mu)$ for \emph{every} $\mu \in \cM$. That is, for any innovation distribution in the class,
\[
    \phi_*(\lambda) \leq \phi(\cF_\lambda^\mu) \leq \phi^*(\lambda) \quad \text{for all } \mu \in \cM.
\]
\end{theorem}

\begin{proof}
Immediate from the definitions~\eqref{eq:phi-star-lower}--\eqref{eq:phi-star-upper}. The infimum and supremum are taken over all $\mu \in \cM$, so any particular $\mu$ yields a value within the bounds.
\end{proof}

\begin{remark}[Interpretation]
The bounds $[\phi_*(\lambda), \phi^*(\lambda)]$ provide \emph{robust inference} that does not depend on correctly specifying $\mu$. If $\phi_*(\lambda) \geq c$, then the surrogate is guaranteed to satisfy $\phi \geq c$ for \emph{all} innovation distributions in $\cM$. Conversely, if $\phi^*(\lambda) < c$, the surrogate fails the threshold under all $\mu \in \cM$.

The cost of this robustness is interval width: $[\phi_*(\lambda), \phi^*(\lambda)]$ will generally be wider than a confidence interval for $\phi(\cF_\lambda^\mu)$ at any single $\mu$. The width quantifies sensitivity to $\mu$-misspecification.
\end{remark}

\subsection{Choosing the class $\cM$}

The choice of $\cM$ determines which innovation distributions are considered ``plausible.'' Common choices include:

\paragraph{Dirichlet family.}
Under Assumption~\ref{ass:finite}, a natural choice is
\[
    \cM_{\text{Dir}}(\alpha_{\min}, \alpha_{\max}) = \{\text{Dirichlet}(\alpha, \ldots, \alpha) : \alpha \in [\alpha_{\min}, \alpha_{\max}]\}.
\]
The uniform Dirichlet ($\alpha = 1$) treats all innovation directions equally; smaller $\alpha$ (e.g., $\alpha < 0.1$) favor sparse innovations (point masses on few observations); larger $\alpha$ (e.g., $\alpha > 10$) concentrate near the uniform distribution. A broad range like $[\alpha_{\min}, \alpha_{\max}] = [0.01, 100]$ provides comprehensive robustness.

\paragraph{Vertex augmentation.}
Point masses $\delta_i$ (putting all weight on observation $i$) represent extreme reweighting. Including these ensures coverage for innovations that drastically upweight single observations:
\[
    \cM_{\text{vertex}} = \{\delta_1, \ldots, \delta_n\}.
\]

\paragraph{Combined class.}
For maximal robustness, combine:
\[
    \cM = \cM_{\text{Dir}}(\alpha_{\min}, \alpha_{\max}) \cup \cM_{\text{vertex}} \cup \{\text{uniform on } \Omega\}.
\]
This covers smooth reweighting (Dirichlet), extreme reweighting (vertices), and no reweighting (uniform).

\subsection{Computation via grid search}

Computing the bounds $[\phi_*(\lambda), \phi^*(\lambda)]$ requires optimizing over $\cM$. For the Dirichlet family, we use grid search:

\begin{enumerate}
    \item \textbf{Construct grid:} Select $\alpha_1 < \alpha_2 < \cdots < \alpha_K$ spanning $[\alpha_{\min}, \alpha_{\max}]$ (e.g., log-spaced with $K = 40$ points).

    \item \textbf{For each} $\alpha_k$, compute $\hat{\phi}(\alpha_k) = \phi(\cF_\lambda^{\text{Dir}(\alpha_k)})$ using the plug-in estimator:
          \begin{enumerate}
              \item Generate $M$ innovations $\tilde{\PP}_m \sim \text{Dirichlet}(\alpha_k, \ldots, \alpha_k)$
              \item Form $\Q_m = (1-\lambda)\hat{\PP}_n + \lambda\tilde{\PP}_m$
              \item Compute $(\hat{\Delta}_S(\Q_m), \hat{\Delta}_Y(\Q_m))$
              \item Compute $\hat{\phi}(\alpha_k)$ as the sample correlation (or other functional)
          \end{enumerate}

    \item \textbf{For each vertex} $i = 1, \ldots, n$ (or a subsample if $n$ is large), compute $\hat{\phi}(\delta_i)$ by setting $\tilde{\PP}_m = \delta_i$ for all $m$.

    \item \textbf{Take extrema:}
          \[
              \hat{\phi}_*(\lambda) = \min\{\hat{\phi}(\alpha_1), \ldots, \hat{\phi}(\alpha_K), \hat{\phi}(\delta_1), \ldots\}
          \]
          \[
              \hat{\phi}^*(\lambda) = \max\{\hat{\phi}(\alpha_1), \ldots, \hat{\phi}(\alpha_K), \hat{\phi}(\delta_1), \ldots\}
          \]
\end{enumerate}

\paragraph{Computational cost.}
Each grid point requires $M$ Monte Carlo draws; with $K$ Dirichlet points and $n$ vertices (or a subsample), total cost is $O((K + n)M)$. For $K = 40$, $n = 1000$ (subsampled to 50 vertices), $M = 1000$: approximately 90,000 treatment effect computations per baseline, taking 5--10 seconds on standard hardware.

\subsection{Implementation via deterministic reweighting}

When computing $\hat{\Delta}_S(\Q_m)$ and $\hat{\Delta}_Y(\Q_m)$ for $\Q_m = (1-\lambda)\hat{\PP}_n + \lambda\tilde{\PP}_m$, we use \emph{deterministic reweighting} of the observed data. For each innovation $\tilde{\PP}_m$, compute observation weights $w_i^{(m)} = q_m(O_i)$ where $q_m$ is the density of $\Q_m$. Then estimate treatment effects via weighted means:
\begin{align*}
    \hat{\Delta}_S(\Q_m) &= \frac{\sum_i w_i^{(m)} A_i S_i}{\sum_i w_i^{(m)} A_i} - \frac{\sum_i w_i^{(m)} (1-A_i) S_i}{\sum_i w_i^{(m)} (1-A_i)}.
\end{align*}
This approach deterministically evaluates the treatment effects under each mixture distribution $\Q_m$, allowing us to explore the space of distributions within $\cF_\lambda$ without introducing additional sampling variability. For constructing confidence intervals on the bounds themselves, standard bootstrap resampling of the original data (recomputing $\hat{\phi}_*(\lambda)$, $\hat{\phi}^*(\lambda)$ on each bootstrap sample) remains appropriate.

\subsection{Asymptotic theory for minimax bounds}

Under the conditions of Proposition~\ref{prop:asymptotic}, each $\hat{\phi}(\mu)$ satisfies
\[
    \sqrt{n}(\hat{\phi}(\mu) - \phi(\cF_\lambda^\mu)) \xrightarrow{d} N(0, \sigma^2_\mu(\lambda))
\]
where $\sigma^2_\mu(\lambda)$ depends on $\mu$ through the gradient of $H$ and the variance matrix $V(\lambda)$.

For the bounds, we have:

\begin{proposition}[Consistency of minimax bounds]\label{prop:minimax-consistency}
Under Assumptions~\ref{ass:bounded}--\ref{ass:nondegeneracy}, as $n \to \infty$ and $M \to \infty$,
\[
    \hat{\phi}_*(\lambda) \xrightarrow{p} \phi_*(\lambda), \quad \hat{\phi}^*(\lambda) \xrightarrow{p} \phi^*(\lambda).
\]
If the grid is sufficiently dense (i.e., the discretization error is $o_p(n^{-1/2})$), then
\[
    \sqrt{n}(\hat{\phi}_*(\lambda) - \phi_*(\lambda)) = O_p(1), \quad \sqrt{n}(\hat{\phi}^*(\lambda) - \phi^*(\lambda)) = O_p(1).
\]
\end{proposition}

\begin{proof}[Proof sketch]
Each $\hat{\phi}(\mu)$ is consistent for $\phi(\cF_\lambda^\mu)$ by Proposition~\ref{prop:asymptotic}. The bounds are continuous functionals (min/max) of the individual estimates. If the grid is dense enough that $\sup_{\mu \in \cM} |\phi(\cF_\lambda^\mu) - \phi(\cF_\lambda^{\mu_{\text{grid}}})| = o(n^{-1/2})$ for some $\mu_{\text{grid}}$ on the grid, then the discretization error is negligible and consistency follows. Standard $M$-estimation results apply to the extremum.
\end{proof}

\paragraph{Bootstrap inference for bounds.}
To quantify uncertainty in the bounds themselves, we can bootstrap:
\begin{enumerate}
    \item For $b = 1, \ldots, B$, resample the data to get $\hat{\PP}_n^{(b)}$
    \item For each bootstrap sample, compute $\hat{\phi}_*^{(b)}(\lambda)$ and $\hat{\phi}^{*(b)}(\lambda)$ via the grid search
    \item Use percentile confidence intervals on the bounds
\end{enumerate}

This is computationally intensive (requires $B \times (K + n) \times M$ evaluations), but provides honest uncertainty quantification on the worst-case bounds.

\subsection{Comparison with standard method}

The standard method (Section~\ref{sec:estimation}) assumes $\mu$ is known (e.g., $\mu = \text{Dirichlet}(1, \ldots, 1)$) and provides a confidence interval $[\ell, u]$ for $\phi(\cF_\lambda^\mu)$. The minimax bounds provide $[\phi_*(\lambda), \phi^*(\lambda)]$ that hold for all $\mu \in \cM$.

\begin{itemize}
    \item \textbf{If} $[\phi_*(\lambda), \phi^*(\lambda)]$ is narrow and $[\ell, u] \supset [\phi_*(\lambda), \phi^*(\lambda)]$: the standard method adequately quantifies $\mu$-uncertainty, and the surrogate quality is robust to $\mu$ choice.

    \item \textbf{If} $[\phi_*(\lambda), \phi^*(\lambda)]$ is wide: surrogate quality is highly sensitive to $\mu$. The standard method may be overconfident if $\mu$ is misspecified.

    \item \textbf{If} $\phi_*(\lambda) \geq c$: the surrogate is guaranteed to meet the threshold for all $\mu \in \cM$, providing strong evidence of robust surrogate quality.
\end{itemize}

\begin{remark}[Width-robustness trade-off]
The width $\phi^*(\lambda) - \phi_*(\lambda)$ quantifies the cost of not knowing $\mu$. In validation studies (Section 5), we observe width ratios of 5--6$\times$ compared to standard confidence intervals: minimax bounds provide guaranteed coverage at the cost of 5--6$\times$ wider intervals. This trade-off is acceptable when $\mu$-misspecification risk is high.
\end{remark}

\subsection{Practical recommendations}

\begin{itemize}
    \item \textbf{Class choice:} Start with $\cM = \cM_{\text{Dir}}(0.01, 100) \cup \cM_{\text{vertex}}$. If bounds are too wide, restrict to $\cM_{\text{Dir}}(0.1, 10)$ with justification.

    \item \textbf{Grid density:} Use $K = 40$ Dirichlet points (log-spaced) for publication. $K = 20$ suffices for exploratory analysis.

    \item \textbf{Vertex subsampling:} For $n > 100$, subsample 50 vertices randomly. This captures extremes without excessive computation.

    \item \textbf{Reporting:} Report both $[\phi_*(\lambda), \phi^*(\lambda)]$ (minimax bounds) and $[\ell, u]$ (standard CI at $\mu = \text{Dirichlet}(1, \ldots, 1)$) for comparison.

    \item \textbf{Sensitivity analysis:} Plot $\hat{\phi}(\alpha)$ vs. $\log(\alpha)$ to visualize sensitivity to $\mu$ across the Dirichlet family. Identify which $\alpha$ values achieve the extrema.
\end{itemize}

\section{Inference for $\varepsilon$-close statements}\label{sec:grid}

The asymptotic theory in Section~\ref{sec:estimation} applies to a single fixed
$\lambda$. In practice, we want to find the largest $\lambda^*$ such that
$\phi(\cF_\lambda) \geq c$ (where $c$ is a chosen threshold, e.g., correlation
$\geq 0.8$). This requires evaluating $\phi(\cF_\lambda)$ over a grid of $\lambda$
values, raising multiple testing concerns.

\subsection{Grid search algorithm}

To find $\lambda^*$:

\begin{enumerate}
    \item \textbf{Choose grid:} Select $\lambda_1 < \lambda_2 < \cdots < \lambda_K$
          (e.g., $\{0.05, 0.10, \ldots, 0.50\}$).

    \item \textbf{For each} $k = 1, \ldots, K$:
          \begin{enumerate}
              \item Fix $\lambda = \lambda_k$
              \item Generate $M$ innovations $\tilde{\PP}_m \sim \mu$
              \item Form $\Q_m = (1-\lambda_k)\hat{\PP}_n + \lambda_k\tilde{\PP}_m$
              \item Compute $(\hat{\Delta}_S(\Q_m), \hat{\Delta}_Y(\Q_m))$
              \item Compute $\hat{\phi}(\lambda_k)$ (e.g., sample correlation)
              \item Bootstrap to get confidence interval $[\ell_k, u_k]$
          \end{enumerate}

    \item \textbf{Find crossing point:} Define
          \[
              \hat{\lambda}^* = \max\{\lambda_k : \ell_k \geq c\}
          \]
          where $\ell_k$ is the lower confidence bound. This ensures
          $\phi(\cF_{\lambda^*}) \geq c$ with confidence $1-\alpha$.
\end{enumerate}

Because $d_{\mathrm{TV}}(\Q, \PP_0) \leq \lambda$ (equation~\eqref{eq:distance-bound}),
we set $\hat{\varepsilon}^* = \hat{\lambda}^*$. The $\varepsilon$-close statement is:
``the surrogate remains good (correlation $\geq c$) as long as future studies are
within total variation distance $\hat{\varepsilon}^*$ of the current study.''

\subsection{Multiple testing adjustment}

Evaluating $\hat{\phi}(\lambda_k) \geq c$ at $K$ values raises multiplicity concerns.
We recommend one of:

\begin{enumerate}
    \item \textbf{Bonferroni correction:} Use significance level $\alpha/K$ for
          each bootstrap confidence interval. Conservative but simple.

    \item \textbf{Simultaneous confidence bands:} Construct bands for
          $\{\phi(\lambda_k)\}_{k=1}^K$ using multiplicity-adjusted bootstrap:
          resample data $B$ times; for each resample, compute $\{\hat{\phi}^{(b)}(\lambda_k)\}_{k=1}^K$;
          take pointwise percentiles adjusted by $\sup_{k} |\hat{\phi}^{(b)}(\lambda_k) - \hat{\phi}(\lambda_k)|$.
          Computationally intensive but controls familywise error exactly.

    \item \textbf{Sequential testing:} If $\phi(\lambda)$ is known or verified to be
          monotone decreasing in $\lambda$, test sequentially from smallest to largest
          $\lambda_k$ until $\hat{\phi}(\lambda_k) < c$. Stops at first crossing.
\end{enumerate}

We recommend option (2) for formal inference, or (1) for simplicity.

\subsection{Practical recommendations}

\begin{itemize}
    \item \textbf{Grid spacing:} Use finer spacing (e.g., 0.025) near the expected
          crossing point; coarser spacing (0.05 or 0.1) elsewhere.
    \item \textbf{Monotonicity:} Check empirically whether $\hat{\phi}(\lambda)$
          is monotone. Under uniform $\mu = \text{Dirichlet}(1,\ldots,1)$,
          correlation typically decreases with $\lambda$, but verify.
    \item \textbf{Computational cost:} Each $\lambda_k$ requires $M$ Monte Carlo
          draws and $B$ bootstrap replicates. Total cost: $O(KMB)$. Parallelize
          over $k$.
\end{itemize}

\section{Simulation study}

\subsection{Design}

The simulation study evaluates:
\begin{enumerate}
    \item \textbf{Finite-sample performance:} Coverage of bootstrap confidence
          intervals for $\phi(\lambda)$ under correctly specified models.

    \item \textbf{Comparison with traditional methods:} Scenarios where
          within-study correlation differs from across-study correlation. DGPs include:
          \begin{itemize}
              \item Strong within-study S-Y correlation but weak treatment effect
                    correlation (spurious surrogate)
              \item Weak within-study correlation but strong treatment effect
                    correlation (good but non-obvious surrogate)
          \end{itemize}

    \item \textbf{Sensitivity to $\mu$:} Compare $\hat{\phi}(\lambda)$ under
          $\mu = \text{Dirichlet}(1,\ldots,1)$ vs. $\mu = \text{Dirichlet}(\alpha\bp_0)$.

    \item \textbf{Grid search performance:} Evaluate multiple testing adjustments
          from Section~\ref{sec:grid}; compare Bonferroni vs. simultaneous bands.
\end{enumerate}

Settings: Sample sizes $n \in \{250, 500, 1000\}$, functionals (correlation,
probability, conditional mean), study types (randomized, observational with
moderate confounding), $B = 500$ bootstrap replicates, $M = 1000$ Monte Carlo draws,
grid spacing $\Delta\lambda = 0.05$.

\subsection{Results}

Validation studies demonstrate that the minimax approach achieves high accuracy in approximating the TV-ball minimax. With sample sizes $n \geq 1000$ and appropriate discretization ($J \in [9, 25]$ bins), approximation error is consistently below $2\%$ across correlation values ranging from $-0.8$ to $0.95$. The method shows excellent convergence properties: error decreases by approximately $75\%$ when sample size increases from $n = 500$ to $n = 4000$, consistent with $O(n^{-1/2})$ or better rates.

The ensemble approach---taking the minimum over multiple discretization schemes---consistently outperforms single-scheme approaches by $10$--$20\%$, validating that different schemes explore different aspects of the treatment effect distribution. RF-based discretizations, which adaptively partition based on estimated treatment effects, perform particularly well when combined with covariate-based schemes in the ensemble.

\section{Theoretical properties of the RF-ensemble approximation}

The minimax bounds $[\phi_*(\lambda), \phi^*(\lambda)]$ provide a principled answer to the question: \emph{what is the worst-case surrogate performance over all studies within total variation distance $\lambda$ of $\PP_0$?} While Section~\ref{sec:minimax} established computation and finite-sample properties, we now address the fundamental question: do these bounds converge to the TV-ball minimax as $n \to \infty$?

\subsection{TV-ball minimax as the target}

For a fixed $\lambda$, define the TV-ball around $\PP_0$ as
\[
    \cB_\lambda(\PP_0) = \{\Q \in \cM_1(\Omega) : d_{\mathrm{TV}}(\Q, \PP_0) \leq \lambda\}.
\]
The \emph{TV-ball minimax} is the worst-case functional value over this ball:
\begin{align}\label{eq:tv-minimax}
    \phi_{\min}^{\mathrm{TV}}(\lambda) = \inf_{\Q \in \cB_\lambda(\PP_0)} \phi(\Q).
\end{align}
This represents an absolute worst-case: no matter which innovation distribution $\mu$ we specify, $\phi(\cF_\lambda^\mu) \geq \phi_{\min}^{\mathrm{TV}}(\lambda)$ because the support of $\cF_\lambda^\mu$ is contained in $\cB_\lambda(\PP_0)$.

The question is: does our discretization-based method approximate $\phi_{\min}^{\mathrm{TV}}(\lambda)$?

\subsection{RF-ensemble approximation theorem}

When covariates are continuous and treatment effects $\tau_S(\bX)$, $\tau_Y(\bX)$ vary smoothly, the finite-support assumption requires discretization. Different discretization schemes explore different aspects of the treatment effect distribution. The following informal theorem characterizes convergence:

\begin{theorem}[RF-Ensemble Approximation, Informal]\label{thm:rf-ensemble}
Suppose treatment effects $\tau_S(\bX)$ and $\tau_Y(\bX)$ are Lipschitz continuous with respect to $\bX$. Let $\cJ = \{J_1, \ldots, J_K\}$ be a collection of discretization schemes with $J_k \to \infty$ as $n \to \infty$. For each scheme $k$, let $\hat{\phi}_{\min}^{(k)}(\lambda)$ be the estimated minimax from Eq.~\eqref{eq:phi-star-lower} using $J_k$ bins. Define the ensemble estimator
\[
    \hat{\phi}_{\min}^{\mathrm{ens}}(\lambda) = \min_{k=1,\ldots,K} \hat{\phi}_{\min}^{(k)}(\lambda).
\]
Under regularity conditions, as $n \to \infty$ and $J_k \to \infty$ with $J_k = o(n)$,
\[
    \hat{\phi}_{\min}^{\mathrm{ens}}(\lambda) \xrightarrow{p} \phi_{\min}^{\mathrm{TV}}(\lambda).
\]
\end{theorem}

\paragraph{Proof strategy (sketch).}
\begin{enumerate}
    \item \textbf{Random forest consistency}: RFs trained on treatment effect estimates $\hat{\tau}(\bX)$ consistently estimate the true functions $\tau(\bX)$ (Wager \& Athey, 2018). The induced partition creates regions where $\tau(\bX)$ is approximately constant.

    \item \textbf{Discretization approximation}: Any distribution $\Q \in \cB_\lambda(\PP_0)$ induces a distribution $G^\Q$ over treatment effect pairs $(\tau_S(\bX), \tau_Y(\bX))$. With $J \to \infty$, discretized distributions $G^{\Q, J}$ approximate $G^\Q$ in Wasserstein distance.

    \item \textbf{Functional continuity}: The functional $\phi$ (e.g., correlation) varies smoothly with the induced distribution $G$. Wasserstein continuity ensures $|\phi(G^{\Q, J}) - \phi(G^\Q)| = O(J^{-\alpha})$ for $\alpha > 0$ depending on smoothness.

    \item \textbf{Ensemble covering}: Using multiple discretization schemes (RF-based, quantile-based, k-means) ensures that for any adversarial $\Q^* \in \cB_\lambda$, at least one scheme produces a partition that captures the heterogeneity induced by $\Q^*$. Taking the minimum over schemes yields $\inf_k \phi(G^{\Q^*, J_k}) \to \phi(G^{\Q^*})$.

    \item \textbf{Estimation consistency}: For each scheme, the plug-in estimator $\hat{\phi}_{\min}^{(k)}$ converges to the discretized minimax. The minimum over finitely many consistent estimators is consistent.
\end{enumerate}

\subsection{Practical implications}

\paragraph{Approximation quality.}
Empirical validation confirms the theorem: with $n = 2000$ and $J \in \{9, 16, 25\}$, approximation error is consistently below $2\%$ across correlation values in $[-0.8, 0.95]$, demonstrating that the RF-ensemble with deterministic reweighting accurately approximates the TV-ball minimax.

\paragraph{Convergence rate.}
Observed error reduction from $n = 500$ ($\sim 10\%$ error) to $n = 4000$ ($\sim 1\%$ error) is consistent with $O(n^{-1/2})$ or better, suggesting the method achieves near-parametric rates despite the nonparametric discretization.

\paragraph{Number of schemes.}
Using $K = 3$--$5$ diverse schemes (e.g., RF-based, age-risk, age-biomarker, risk-biomarker, k-means) provides substantial improvement ($10$--$20\%$) over single-scheme approaches while keeping computation tractable. Additional schemes beyond $K = 5$ show diminishing returns.

\paragraph{Choice of $J$.}
For each scheme, using $J \in \{9, 16, 25\}$ and taking the minimum balances discretization error (decreases with $J$) and estimation noise (increases with $J$). Adaptive choice based on effective sample size per bin (target: $\geq 20$ observations/bin) works well in practice.

\subsection{Limitations and extensions}

\paragraph{Known limitations.}
The method requires adequate treatment effect heterogeneity: when treatment effects are near-constant across $\bX$, all future studies yield similar $(\Delta_S, \Delta_Y)$ pairs regardless of covariate distribution shifts, making surrogate evaluation trivial (but not problematic). For weak correlations ($\rho < 0.5$) with many latent subgroups ($K \geq 50$), discretization with moderate $J$ may under-represent heterogeneity, leading to conservative (wider) bounds.

\paragraph{Continuous approximation.}
An alternative to discretization is kernel-based estimation directly on continuous $\bX$. However, this requires bandwidth selection and scales poorly to high-dimensional $\bX$. The discretization approach is more robust and computationally efficient.

\paragraph{Formal proof.}
A rigorous proof of Theorem~\ref{thm:rf-ensemble} requires formalizing the covering property of the ensemble and establishing uniform convergence rates. This is left for future work but is strongly supported by the empirical validation.

% To operationalize this, we assume that $\PP_0 \sim \text{Dirichlet}(\balpha)$ for $\balpha = (\alpha_1, ..., \alpha_n), \alpha_k >0$, which results in a posterior distribution of $\cF | \PP_0 \sim \text{Dirichlet}(\balpha + 1).$ Choosing $\balpha$ encodes how similar new studies are expected to be to the current study on $\PP_0$. As $\balpha \rightarrow \infty$, new studies will be identical to the current study. 

% To better understand how choices of $\balpha$ impact the posterior distribution, consider the toy examples in Figure \ref{fig:alpha-toy-example}. When the $\balpha$ is uniform, the distribution of future studies is centered on $\PP_0$, with the magnitude of $\balpha$ determining whether future studies are mostly similar to the current study (top right in Figure \ref{fig:alpha-toy-example}, $\balpha = (5, 5, 5)$) or more diffuse (top left in the figure). When $\balpha$ is not uniform, the posterior distribution may be systematically different from the current study, primarily including studies with certain types of individuals over-represented compared to $\PP_0$ and others under-represented (bottom row of the figure). 

% \begin{figure}[htbp]
%     \centering
%     \includegraphics[width=0.8\textwidth]{dirichlet_contours_clean.png}
%     \caption{Toy examples of $\balpha$ impact on posterior distribution $\cF | \PP_0$ when $n = 3$. Contours indicate the density of the posterior, with lighter colors indicating higher density. Areas near the middle of the triangle correspond to similar studies to the current study, while areas near the edges correspond to studies with distributions that differ more from $\PP_0$.}
%     \label{fig:alpha-toy-example}
% \end{figure}

% \subsection{Procedure for a pre-specified $\balpha$}

% Depending on the scientific context, researchers may choose a specific $\balpha$ that describes their knowledge of future studies. For example, if the current study includes both young and old patients but most future studies will be focused primarily on older patients, they may choose $\balpha$ to be weighted more toward older patients. Or $\balpha$ may be chosen to be uniform such that $\alpha_k = c$ for each $k$ and some $c > 0$, where $c$ may be chosen to control the spread in the posterior distribution. 

% Let's consider some examples. We will stick with the example in \ref{tab:binary_combinations} for the moment, and we will select $\cF$ to be Dirichlet($\bpi$), where $\bpi = (\pi_{0000}, \pi_{0001}, ..., \pi_{1111})$. If we select an improper/uninformative prior on $\cF$, where the prior parameters don't contribute anything to the resulting posterior, then we obtain $\cF|\P_0 \sim \text{Dirichlet}(K\bp)$, where $\bp = (p_{0000}, ..., p_{1111})$, the observed probability masses, and $K$ is the tuning parameter controlling the spread of the distribution. 

% \subsection{Using correlation as the measure}

% If what we want is to estimate (I realize this doesn't fit exactly into the format of \eqref{posterior}, but let's worry about that later)
% \begin{align*}
%     \cor_{\cF | \P_0}\{\Delta_y(\Q), \Delta_S(\Q)\}
% \end{align*}
% then we can recover this in closed form, and the value doesn't depend on the tuning parameter $K$. In particular, we have that
% \begin{align*}
%     \cor_{\cF | \P_0}\{\Delta_y(\Q), \Delta_S(\Q)\} = \frac{\ba\trans\Sigma\bb}{\sqrt{\ba\trans\Sigma\ba \bb\trans\Sigma\bb}}
% \end{align*}
% where $\ba = (0, 1, -1, 0, 0, 1, -1, 0, 0, 1, -1, 0, 0, 1, -1, 0)$ such that $\ba\trans\bp = \Delta_Y$ and $\bb = (0, 0, 0, 0, 1, 1, 1, 1, -1, -1, -1, -1, 0, 0, 0, 0)$ such that $\bb\trans\bp = \Delta_S$, and 
% \begin{align*}
%     \Sigma = \Cov(\bp) = \inv{(K+1)} (\text{diag}(\bp) - \bp\bp\trans).
% \end{align*}

% We will put off for a moment considering other effect measures beyond correlation to consider how this can be adapted to settings on real data.

% \subsection{Extending this approach to real data}

% In real data, we don't have the potential outcomes. It might at first seem sensible to use the empirical means and the distribution $\bptilde = 1/n$ for each observation. Now we have $\Deltahat_Y = \ba\trans\bptilde$ where $\ba = Y\{A/\E(A) - (1-A)/\E(1-A)\}$ and $\Deltahat_S = \bb\trans\bp$ where $\bb = S\{A/\E(A) - (1-A)/\E(1-A)\}.$ However, computing
% \begin{align*}
%     \cor_{\cF | \P_0}\{\Deltahat_y(\Q), \Deltahat_S(\Q)\} = \frac{\ba\trans\Sigma\bb}{\sqrt{\ba\trans\Sigma\ba \bb\trans\Sigma\bb}}
% \end{align*}
% will be inflated by the correlation induced by $A$, and it doesn't immediately seem apparent how to avoid this.

% % Instead, we can acknowledge that we can't observe $\bp$ directly, but we can bound it given the data we observe. Specifically, we require that
% % \begin{align}
% %     \Omega\bp = \obar\label{constraint}
% % \end{align}
% % where $\obar$ is a vector containing values of $\Phat_n(A = a, S = s, Y = y)$, and $\Omega_{8\times16}$ is a binary constraint matrix that indicates which of the 16 counterfactual types is consistent with each of the eight observed $(A, S, Y)$ tuples and ensures that the distribution of the counterfactuals is consistent with the empirical distribution $\obar$. We may then compute $\cor_{\cF | \P_0}\{\Deltahat_y(\Q), \Deltahat_S(\Q)\}$ under the constraint \eqref{constraint}.

% Instead, we can write down the likelihood for the observed data in terms of $\bp$, which is
% \begin{align*}
%     \cL(\bp) = \prod_{s = 0}^1\prod_{y=0}^1 \left\{(1-\pi)\sum_{s' = 0}^1\sum_{y'=0}^1 p_{ss'yy'}\right\}^{n_{0sy}}\left\{\pi\sum_{s' = 0}^1\sum_{y'=0}^1 p_{s'sy'y}\right\}^{n_{1sy}}
% \end{align*}
% where $\pi= \P_0(A = 1)$ and $n_{asy} = \sum_{i=1}^n\cI\{S = s, Y = y, A = a\}$.
% Recalling that $\bp \sim$ Dirichlet($\bpi$), we put a prior on the entries in $\bpi$
% \begin{align*}
%     \pi_{ss'yy'} \sim \text{Gamma}(\alpha_{ss'yy'}, \beta_{ss'yy'}),
% \end{align*}
% and we compute the posterior distribution of $\bpi$.

\bibliography{refs}
\bibliographystyle{plain}
\end{document}
